% Section 19.2, translated from sv/chapters/19/practice_exam.tex. The answers are in
% en/back/facit/19.tex.

\section{Practice exam}
\label{c19:sec:ovningstentamen}
The practice exam has the same aids, points and grade thresholds as the exam. Like the exam, it is
divided into two parts: Part~I (Chapters~\ref{c1:ch}--\ref{c10:ch}) and Part~II
(Chapters~\ref{c12:ch}--\ref{c17:ch}).

\begin{aside}
\textbf{How to work with the practice exam.} Sit it on your own and under the same conditions as
the exam: with only the permitted aids, without a mobile phone and without looking at the answer
key. Work each question as far as you can before you move on.

\textbf{Answers and worked solutions.} The answer to every question is in the answer key at the
back of the book, and its page is given next to the question. An answer shows whether you got it
right, but not how. Full worked solutions, in Swedish, are published in the course repository, in
\file{exam/practice_exam_solutions.md}. Always try it yourself first.
\end{aside}

\begin{booktable}{@{}p{7.2em}L@{}}
  \thead{Aids} & Writing materials and a calculator of any kind. A formula sheet on one A4 sheet
    (both sides), handwritten or typed, and one handwritten A4 sheet (both sides) with notes of
    your own choice. \\
  \thead{Points} & $25$ points in total. On top of that, up to $1.0$ bonus point can be earned
    from tests~1--2. \\
  \thead{Grades} & IG $< 10$ points, G $\geq 10$ points, VG $\geq 17.5$ points. \\
  \thead{Solutions} & Every question requires a solution and a justification with a stated
    answer, including the unit where relevant. In other words, you must explain how you arrived
    at your results. Good luck! \\
\end{booktable}

\begin{aside}
\note[Note] Mobile phones may not be used while the exam is in progress and must be left in the
place indicated.
\end{aside}

\prov{ot}

% ------------------------------------------------------------------------------------------
\uppgiftsdel{Part I}{Arithmetic, algebra, equations, trigonometry and logarithms}

\provuppgift{1}{1.0}
The resistance of four resistors $R_1$, $R_2$, $R_3$ and $R_4$ connected in parallel can be
calculated with the formula
\[
  \frac{1}{R} = \frac{1}{R_1} + \frac{1}{R_2} + \frac{1}{R_3} + \frac{1}{R_4}.
\]
Calculate the parallel resistance $R$ if $R_1 = R_2 = 2.2\,\text{k}\Omega$ and
$R_3 = R_4 = 10\,\text{k}\Omega$. Give the answer in k$\Omega$ to one significant figure.

\provuppgift{2}{1.0}
Simplify the following expression as far as possible, without regard to invalid values of $c$:
\[
  \frac{\dfrac{c^8 - 16c^4}{c^2}}{c(c^2 - 4)}
\]

\provuppgift{3}{1.0}
Solve the system of equations
\[
  \begin{cases}
    4y = -6x - 6 \\
    3y + 6x + 9 = 0
  \end{cases}
\]

\provuppgift{4}{1.0}
\begin{parts}
  \item The equations in \hyperref[ot:u:3]{question~3} are examples of the equation of a straight
    line. Give the first equation's slope $k$ and its intercept $m$.
  \item The equation
    \[
      \sin(v) = -\frac{\sqrt{3}}{2}
    \]
    has a solution in the interval between $180^{\circ}$ and $270^{\circ}$. Find this solution and
    give it in degrees.
\end{parts}

\provuppgift{5}{2.0}
Solve the equation
\[
  (x + 2)^2 + 2(x - 2)^2 = 45 - 6x.
\]

\provuppgift{6}{1.0}
You have two vectors $u = (3, 5)$ and $v = (2, -4)$.
\begin{parts}
  \item Draw the vectors in a coordinate system.
  \item Calculate the magnitudes of the vectors, $\abs{u}$ and $\abs{v}$.
  \item Calculate the angles of the vectors, $v_u$ and $v_v$.
  \item Calculate a third vector $w = u - 2v$.
\end{parts}

\provuppgift{7}{1.0}
House prices rise over time because of inflation and market growth. The price $f(x)$ of a home
can be described by the exponential function
\[
  f(x) = C \cdot a^x,
\]
where $C$ is the original price, $a$ is the annual growth factor and $x$ is the number of years
that have passed since the original price was set.

If, for example, the annual price increase is $5\,\%$ over a certain period, a home that costs
$3\,000\,000\,\text{kr}$ today is worth $f(x)$~kr after $x$ years, where
\[
  f(x) = 3 \cdot 10^6 \cdot 1.05^x.
\]
Assume that the original price of a home is $7\,500\,000\,\text{kr}$ and that the annual price
increase is $2.5\,\%$.
\begin{parts}
  \item Give a function that describes the home's price $f(x)$ after $x$ years.
  \item Calculate the home's value after $5$ years.
  \item Find the function's domain and range for a period of $0$--$30$ years.
  \item After how many years has the home's value doubled?
\end{parts}

\provuppgift{8}{1.0}
Simplify the following rational expression and state which values of $x$ it is not defined for:
\[
  f(x) = \frac{3(x^2 - x - 12)}{x^2 - 9}
\]

\provuppgift{9}{1.0}
When a capacitor discharges through a resistor, the voltage across the capacitor decreases
exponentially with time. The voltage $u(t)$ can be described by the function
\[
  u(t) = U_0 e^{-t/(RC)},
\]
where
\begin{itemize}
  \item $u(t)$ is the voltage across the capacitor at time $t$,
  \item $U_0$ is the initial voltage in volts,
  \item $RC$ is the circuit's time constant in seconds,
  \item $t$ is the time in seconds.
\end{itemize}
A capacitor with the capacitance $680\,\upmu\text{F}$ is connected to a $10\,\text{k}\Omega$
resistor. It is initially charged to $10\,\text{V}$ and starts to discharge at time $t = 0$.
\begin{parts}
  \item Calculate the voltage after three seconds.
  \item Find the function's domain and range for the time interval $0$--$10$ seconds.
  \item Calculate how long it takes for the voltage to fall to $40\,\%$ of its original value.
\end{parts}

\provuppgift{10}{1.5}
An AC voltage has the amplitude $5\,\text{V}$, the frequency $100\,\text{Hz}$ and the phase
$90^{\circ}$.
\begin{parts}
  \item Find the equation $u(t)$ of the AC voltage. Give the phase in radians.
  \item Draw the sine wave of the AC voltage over one period $T$.
\end{parts}

\provuppgift{11}{1.0}
An audio amplifier has a gain of $54\,\text{dB}$. What voltage amplification factor does that
correspond to?

% ------------------------------------------------------------------------------------------
\uppgiftsdel{Part II}{Derivatives, integrals and complex numbers}

\provuppgift{12}{1.0}
Consider the function
\[
  f(x) = x^2 - 4x + 3.
\]
\begin{parts}
  \item Differentiate the function $f(x)$ and find the expression for $f'(x)$.
  \item Find where the function is stationary, i.e.\ solve $f'(x) = 0$.
  \item Use the second derivative to decide whether the point is a maximum or a minimum.
  \item Calculate the smallest or largest value of the function, i.e.\ find $f(x)$ at the point
    where the function is stationary.
  \item Draw the graph of $f(x)$ for the interval $0 \leq x \leq 5$. Mark the turning point and
    any points where the graph crosses the axes.
\end{parts}

\provuppgift{13}{1.0}
The charging of a capacitor can be described by the formula
\[
  u(t) = 10\bigl(1 - e^{-0.5t}\bigr),
\]
where $u(t)$ is the voltage drop across the capacitor at time $t$ and $t$ is the time in seconds.
\begin{parts}
  \item Differentiate the function and find the expression for $u'(t)$.
  \item Calculate $u'(t)$ at time $t = 2\,\text{s}$, i.e.\ calculate $u'(2)$.
  \item Decide whether the rate of increase of the voltage is increasing or decreasing at time
    $t = 2\,\text{s}$, i.e.\ calculate $u''(2)$.
\end{parts}

\provuppgift{14}{1.5}
Differentiate the following functions.
\begin{delar}(2)
  \task $f(x) = e^{2x} \cdot \ln x^2$
  \task $f(x) = \dfrac{3x^2}{\ln x}$
  \task $u(t) = 5\cos\left(100\pi t + \dfrac{\pi}{4}\right)\,\text{V}$
\end{delar}

\provuppgift{15}{1.0}
The charge $q(t)$ that passes through a conductor during the time interval $(0, t)$ can be
calculated with
\[
  q(t) = \int_0^t i(t)\,dt = \bigl[I(t)\bigr]_0^t,
\]
where
\begin{itemize}
  \item $q(t)$ is the charge through the conductor in coulombs (C),
  \item $i(t)$ is the current through the conductor in amperes (A),
  \item $I(t)$ is the antiderivative (indefinite integral) of $i(t)$.
\end{itemize}
The current $i(t)$ in a capacitor is given by
\[
  i(t) = 0.5e^{-0.1t},
\]
where $t$ is the time in seconds. The capacitor already has the charge $q_0 = 1\,\text{C}$ at the
start.
\begin{parts}
  \item Find an expression for the antiderivative $I(t) = \int i(t)\,dt$, i.e.\ integrate
    without any limits.
  \item Find an expression for the charge $q(t)$ of the capacitor by integrating over the interval
    $[0, t]$.
  \item Find a formula for the total charge $q_{\text{tot}}(t)$ in the capacitor, including the
    initial charge $q_0$.
  \item How much total charge (including the initial charge) has flowed into the capacitor after
    $4$ seconds?
\end{parts}

\provuppgift{16}{2.0}
A circuit is supplied with the voltage $U = 3 - j4\,\text{V}$. The circuit has the impedance
$Z = 2 + j2\,\Omega$. Calculate the current $I$ through the circuit with “Ohm's law”, that is,
Ohm's law adapted to complex numbers:
\[
  I = \frac{U}{Z}
\]
Give the current in both rectangular and polar form.

\provuppgift{17}{2.0}
A current $i(t)$ in an AC circuit can be represented by a phasor $I$, which in rectangular form
is written
\[
  I = 10 - j5\,\text{mA}.
\]
\begin{parts}
  \item Draw the phasor $I$ in the complex plane (the $x$-axis is the real part, the $y$-axis the
    imaginary part).
  \item Express the phasor $I$ in Euler form, i.e.\ find the magnitude $\abs{I}$ and the phase
    angle $\delta$ such that $I = \abs{I}e^{j\delta}$.
  \item Assume that the frequency of the current is $f = 10\,\text{Hz}$. Find the angular
    frequency $w$.
  \item Write $i(t)$ as a function of time $i(t) = \abs{I} \cdot e^{j(wt + \delta)}$.
\end{parts}

\provuppgift{18}{2.0}
You have the vectors $a = (2, 1)$, $b = (3, -4)$ and $c = (-1, 3)$. In the parts below, each
vector $(x, y)$ is to be read as a complex number $z = x + jy$.
\begin{parts}
  \item Write the vectors in complex rectangular form.
  \item Draw the vectors in the complex plane (the $x$-axis is the real part, the $y$-axis the
    imaginary part).
  \item Find the lengths of the vectors, i.e.\ the magnitude of each number.
  \item Find the length (magnitude) of $2a - 3c$, i.e.\ $\abs{2a - 3c}$.
  \item Find the angles of the vectors.
  \item Find a vector $d$ with length $7$ that points in the opposite direction to $a$.
\end{parts}

\provuppgift{19}{2.0}
In a series connection of three components, the AC voltage across each component is measured as
\begin{align*}
  u_1(t) &= 2\sin(wt + 25^{\circ})\,\text{V}, \\
  u_2(t) &= 5\sin(wt - 45^{\circ})\,\text{V}, \\
  u_3(t) &= 1.5\sin(wt + 36^{\circ})\,\text{V}.
\end{align*}
The total voltage in the circuit is
\[
  u_{\text{tot}}(t) = u_1(t) + u_2(t) + u_3(t).
\]
\begin{parts}
  \item Rewrite the voltages $u_1(t)$, $u_2(t)$ and $u_3(t)$ as phasors $U_1$, $U_2$ and $U_3$ in
    complex rectangular form.
  \item Calculate the phasor sum $U_{\text{tot}} = U_1 + U_2 + U_3$.
  \item Draw the phasors in the complex plane (the $x$-axis is the real part, the $y$-axis the
    imaginary part).
  \item Convert the result back to a sinusoidal voltage in the time domain, in the form
    $u_{\text{tot}}(t) = \abs{U_{\text{tot}}} \sin(wt + \delta_{\text{tot}})$.
\end{parts}
